\section{Used Hardware}
\label{sec:used_hardware}

The development and evaluation of the Agentic Framework spanned two distinct hardware environments, reflecting the shift from local prototyping to scalable production deployment.

\subsection{Local Prototyping Environment}
Initial development and unit testing were conducted on a workstation equipped with an \textbf{NVIDIA RTX A4000} GPU, featuring 16~GB of dedicated VRAM (GDDR6 with ECC), operating with driver version 575.57.08 and CUDA version 12.9.

The 16~GB VRAM buffer allowed us to deploy quantized Large Language Models (such as Mistral-7B and Qwen-7B via \texttt{Ollama}) directly on the developer's machine. Keeping execution strictly local reduced latency and avoided the privacy risks of cloud-based inference APIs. Local GPU acceleration also enabled rapid iterative testing of complex LangGraph agentic loops. Running these tests on a standard CPU-based environment would have been prohibitively slow. However, this setup proved insufficient for concurrent multi-model access due to severe VRAM limits.

\subsection{High Performance Platform (HPP) Cluster}
To handle production workloads and multi-user stress testing, we migrated the core inference engine to the PoliTo \textbf{High Performance Platform (HPP)}. The HPP cluster gave us access to powerful GPUs with \textbf{40~GB} of VRAM capacity.

This expanded hardware infrastructure was necessary to deploy the \textbf{NVIDIA Triton Inference Server}. The 40~GB capacity allowed multiple heavyweight models (such as Mistral and \texttt{gpt-oss-20b}) to remain resident in memory simultaneously without triggering destructive load/unload cycles. By delegating LLM execution over the network to the HPP hardware, we freed the local workstation CPU and RAM entirely for LangGraph orchestration logic and intensive memory-bound fine-tuning tasks.
